\textbf{Results.} In these runs, what the auditor was told the deliverable is changed its flag rate
substantially (Figure~\ref{fig:ai4s}a). In the registered design the task
opened with the whole multi-step problem, and the shipped auditor flagged 87.7\% [72.7, 98.6] of
defective instances and 68.7\% [56.6, 79.7] of correct ones at eight readings. Many of its
objections to correct code were that one step did not implement the whole problem. A second
registered arm, run after the first and changing only the task text, named the step's function as
the deliverable and gave the problem description as context. On the same program bytes, flags on
correct code fell to 36.0\% [26.4, 45.6], a paired change of $-32.7$ points [$-42.8$, $-22.4$], and
recall fell to 72.8\% [59.3, 87.3], $-14.8$ points [$-28.6$, $-2.7$]. The generator's own model flagged
less of both strata in both arms, 46.9 points less of the defective instances than the shipped
auditor (exact McNemar $p = 1.6\times10^{-9}$ in each arm). With no tests shown, the shipped auditor
discriminates, but modestly: at one reading, 61.7\% of defective against 19.8\% of correct instances.

On reading a sample of its flags on correct code, some point at real divergences between the code
and the step's text or its own documentation, and some at contradictions inside the benchmark's own
instructions or at tests too weak to tell outputs apart. Few defects went unflagged by any of the
sixteen readings: 7 in the first arm and 19 in the second. Most of them the two raters called
undetermined by the prose, 4 of 7 and 13 of 19 (for example an entropy's logarithm base, or the
direction of a minimum-image vector), but the registered test that this share exceeds the general-code
residual's 77.2\% is killed in both arms by its own interval rule. These labels describe the prose,
not why a defect was missed.
